\chapter{Metrik Uzaylar}

Metrik uzaylar, mesafe fonksiyonuna dayanan topolojik yapılardır. Her metrik uzay
doğal olarak bir topolojik uzay oluşturur; bu yapı ``indüklenen topoloji'' olarak adlandırılır.

%-------------------------------------------------------------------
\section{Konu}

\subsection{Metrik Aksiyomları}

\begin{definition}[Metrik Uzay]
$M$ kümesi ve $d: M \times M \to [0,\infty)$ verilsin.
$(M, d)$ bir \textbf{metrik uzay}, $d$ ise \textbf{metrik} ise:
\begin{enumerate}
  \item[(M1)] \textbf{Özdeşlik:} $d(x,y) = 0 \Leftrightarrow x = y$
  \item[(M2)] \textbf{Simetri:} $d(x,y) = d(y,x)$
  \item[(M3)] \textbf{Üçgen:} $d(x,z) \leq d(x,y) + d(y,z)$
\end{enumerate}
\end{definition}

\subsection{Temel Kavramlar}

\begin{itemize}
  \item \textbf{Açık top:} $B(x, r) = \{y \in M : d(x,y) < r\}$
  \item \textbf{Kapalı top:} $\bar{B}(x, r) = \{y \in M : d(x,y) \leq r\}$
  \item \textbf{İndüklenen topoloji:} $\tau_d = \{U : \forall x \in U,\,\exists\varepsilon>0,\,B(x,\varepsilon) \subseteq U\}$
  \item \textbf{Çap:} $\mathrm{diam}(A) = \sup\{d(x,y) : x,y \in A\}$
\end{itemize}

\subsection{Standart Metrikler}

\begin{table}[h]
\centering
\begin{tabular}{lll}
\toprule
\textbf{Metrik} & \textbf{Tanım} & \textbf{Uzay} \\
\midrule
Öklid & $d(x,y) = |x-y|$ & $\mathbb{R}$ \\
Ayrık & $d(x,y) = 0$ veya $1$ & Herhangi $X$ \\
$\ell^\infty$ (max) & $\max_i d_i(x_i,y_i)$ & Çarpım uzayı \\
Sınırlı (capped) & $d'(x,y) = \min(d(x,y),1)$ & Eşdeğer topoloji \\
\bottomrule
\end{tabular}
\caption{Standart metrikler}
\end{table}

%-------------------------------------------------------------------
\section{Teoremler}

\begin{theorem}
Her metrik uzay Hausdorff (T2) ve $1^{\text{st}}$ sayılabilirdir.
\begin{proof}[Kanıt İskeleti]
$B(x, d(x,y)/2)$ ve $B(y, d(x,y)/2)$ ayrık komşuluklardır; $\{B(x,1/n)\}$ yerel bazdır.
\end{proof}
\end{theorem}

\begin{theorem}[Sınırlı Metrik Eşdeğerliği]
$d'(x,y) = \min(d(x,y),1)$ ile $d$ aynı topolojiyi indükler.
\end{theorem}

\begin{theorem}
$\mathbb{R}^n$ üzerindeki max, sum ve Öklid metrikleri topolojik olarak eşdeğerdir.
\end{theorem}

%-------------------------------------------------------------------
\section{Algoritmalar}

\subsection{validate\_metric — $O(|X|^3)$}

\begin{algorithm}
\caption{ValidateMetric$(X, d)$}
\begin{algorithmic}[1]
\For{each $x \in X$}
    \If{$d(x,x) \neq 0$} \Return False \EndIf
    \For{each $y \neq x$}
        \If{$d(x,y) = 0$} \Return False \EndIf
    \EndFor
\EndFor
\For{each $x,y \in X$}
    \If{$d(x,y) \neq d(y,x)$} \Return False \EndIf
\EndFor
\For{each $x,y,z \in X$}
    \If{$d(x,z) > d(x,y) + d(y,z)$} \Return False \EndIf
\EndFor
\State \Return True
\end{algorithmic}
\end{algorithm}

Karmaşıklık: $O(|X|^3)$ — üçgen eşitsizliği dominates.

\subsection{open\_ball — $O(|X|)$}

\begin{algorithm}
\caption{OpenBall$(x, r, X, d)$}
\begin{algorithmic}[1]
\State \Return $\{y \in X : d(x,y) < r\}$
\end{algorithmic}
\end{algorithm}

%-------------------------------------------------------------------
\section{pytop API}

\begin{lstlisting}[language=Python]
from pytop.metric_spaces import (
    FiniteMetricSpace,
    SymbolicMetricSpace,
    validate_metric,
)
from pytop import (
    open_ball,
    closed_ball,
    distance_to_subset,
    diameter_of_subset,
    is_bounded_subset,
    capped_metric,
)
\end{lstlisting}

\texttt{FiniteMetricSpace(carrier=tuple, distance=dict)} — mesafe sözlüğü \texttt{(a,b): float}.

\texttt{open\_ball(fms, point, radius)} $\to$ \texttt{set} döner (Result değil).

\texttt{capped\_metric(distance\_fn, cap=1.0)} $\to$ \texttt{Callable} döner.

%-------------------------------------------------------------------
\section{Örnekler}

\subsection*{Örnek 5.1 — Ayrık Metrik Uzay}

\begin{lstlisting}[language=Python]
points = [1, 2, 3, 4]
dist_discrete = {(a, b): (0 if a == b else 1)
                 for a in points for b in points}
fms = FiniteMetricSpace(carrier=tuple(points), distance=dist_discrete)
print("carrier:", fms.carrier)
print("d(1,2):", fms.distance[(1,2)])
\end{lstlisting}

\begin{verbatim}
carrier: (1, 2, 3, 4)
d(1,2): 1
\end{verbatim}

\subsection*{Örnek 5.3 — open\_ball}

\begin{lstlisting}[language=Python]
print("B(1, 0.5) =", open_ball(fms, 1, 0.5))
print("B(1, 1.5) =", open_ball(fms, 1, 1.5))
\end{lstlisting}

\begin{verbatim}
B(1, 0.5) = {1}
B(1, 1.5) = {1, 2, 3, 4}
\end{verbatim}

\subsection*{Örnek 5.6 — capped\_metric}

\begin{lstlisting}[language=Python]
cap_fn = capped_metric(fms_line.distance, cap=1.0)
print("d(0,3) capped:", cap_fn(0, 3))
print("d(0,1) capped:", cap_fn(0, 1))
\end{lstlisting}

\begin{verbatim}
d(0,3) capped: 1.0
d(0,1) capped: 1.0
\end{verbatim}

%-------------------------------------------------------------------
\section{Alıştırmalar}

\subsection*{Kodlama}

\begin{enumerate}[label=K\arabic*.]
\item $\{A,B,C,D\}$ üzerinde bir metrik tanımlayın ve \texttt{validate\_metric} ile
      doğrulayın. Üçgen eşitsizliğini ihlal eden bir örnek de deneyin.
\item Öklid metrikli $\{0,\ldots,9\}$ üzerinde $B(5,3.0)$ ve $\bar{B}(5,3.0)$
      hesaplayın.
\item \texttt{normalized\_metric} kullanarak bir metriği $[0,1]$'e normalleştirin.
\end{enumerate}

\subsection*{Teori}

\begin{enumerate}[label=T\arabic*.]
\item Her metrik uzayın Hausdorff olduğunu ispatlayın.
\item $d' = \min(d,1)$ ve $d$ aynı topolojiyi indüklediğini gösterin.
\end{enumerate}
